%%
%% This is file `sample-sigconf.tex',
%% generated with the docstrip utility.
%%
%% The original source files were:
%%
%% samples.dtx  (with options: `all,proceedings,bibtex,sigconf')
%% 
%% IMPORTANT NOTICE:
%% 
%% For the copyright see the source file.
%% 
%% Any modified versions of this file must be renamed
%% with new filenames distinct from sample-sigconf.tex.
%% 
%% For distribution of the original source see the terms
%% for copying and modification in the file samples.dtx.
%% 
%% This generated file may be distributed as long as the
%% original source files, as listed above, are part of the
%% same distribution. (The sources need not necessarily be
%% in the same archive or directory.)
%%
%%
%% Commands for TeXCount
%TC:macro \cite [option:text,text]
%TC:macro \citep [option:text,text]
%TC:macro \citet [option:text,text]
%TC:envir table 0 1
%TC:envir table* 0 1
%TC:envir tabular [ignore] word
%TC:envir displaymath 0 word
%TC:envir math 0 word
%TC:envir comment 0 0
%%
%% The first command in your LaTeX source must be the \documentclass
%% command.
%%
%% For submission and review of your manuscript please change the
%% command to \documentclass[manuscript, screen, review]{acmart}.
%%
%% When submitting camera ready or to TAPS, please change the command
%% to \documentclass[sigconf]{acmart} or whichever template is required
%% for your publication.
%%
%%
%\documentclass[sigconf]{acmart}
\documentclass[manuscript, screen]{acmart}
%%
%% \BibTeX command to typeset BibTeX logo in the docs
\AtBeginDocument{%
  \providecommand\BibTeX{{%
    Bib\TeX}}}

%% Rights management information.  This information is sent to you
%% when you complete the rights form.  These commands have SAMPLE
%% values in them; it is your responsibility as an author to replace
%% the commands and values with those provided to you when you
%% complete the rights form.
\setcopyright{acmlicensed}
\copyrightyear{2026}
\acmYear{2026}
\acmDOI{XXXXXXX.XXXXXXX}
%% These commands are for a PROCEEDINGS abstract or paper.
\acmConference[Conference acronym 'XX]{Make sure to enter the correct
  conference title from your rights confirmation email}{June 03--05,
  2018}{Woodstock, NY}
%%
%%  Uncomment \acmBooktitle if the title of the proceedings is different
%%  from ``Proceedings of ...''!
%%
%%\acmBooktitle{Woodstock '18: ACM Symposium on Neural Gaze Detection,
%%  June 03--05, 2018, Woodstock, NY}
\acmISBN{978-1-4503-XXXX-X/2018/06}


%%
%% Submission ID.
%% Use this when submitting an article to a sponsored event. You'll
%% receive a unique submission ID from the organizers
%% of the event, and this ID should be used as the parameter to this command.
%%\acmSubmissionID{123-A56-BU3}

%%
%% For managing citations, it is recommended to use bibliography
%% files in BibTeX format.
%%
%% You can then either use BibTeX with the ACM-Reference-Format style,
%% or BibLaTeX with the acmnumeric or acmauthoryear sytles, that include
%% support for advanced citation of software artefact from the
%% biblatex-software package, also separately available on CTAN.
%%
%% Look at the sample-*-biblatex.tex files for templates showcasing
%% the biblatex styles.
%%

%%
%% The majority of ACM publications use numbered citations and
%% references.  The command \citestyle{authoryear} switches to the
%% "author year" style.
%%
%% If you are preparing content for an event
%% sponsored by ACM SIGGRAPH, you must use the "author year" style of
%% citations and references.
%% Uncommenting
%% the next command will enable that style.
%%\citestyle{acmauthoryear}


%%
%% end of the preamble, start of the body of the document source.
\begin{document}

%%
%% The "title" command has an optional parameter,
%% allowing the author to define a "short title" to be used in page headers.
\title{Synergistic Carbon Practice Allocation and Heterogeneous Coalition Strategies for Smallholder Farmers: A Cooperative Game Theory Perspective}

%%
%% The "author" command and its associated commands are used to define
%% the authors and their affiliations.
%% Of note is the shared affiliation of the first two authors, and the
%% "authornote" and "authornotemark" commands
%% used to denote shared contribution to the research.

\author{Thandava Sunkara}
\affiliation{%
  \institution{Indian Institute of Science}
  \city{Bengaluru}
  \country{India}}
\email{thandavas@iisc.ac.in}

\author{Geetha Charan}
\affiliation{%
  \institution{Indian Institute of Science}
  \city{Bengaluru}
  \country{India}}
\email{geethacharan@iisc.ac.in}

\author{Rohit Suresh}
\affiliation{%
  \institution{Indian Institute of Science}
  \city{Bengaluru}
  \country{India}}
\email{rohitpsuresh@iisc.ac.in}

\author{Priyanka V}
\affiliation{%
  \institution{Indian Institute of Science}
  \city{Bengaluru}
  \country{India}}
\email{priyankav@iisc.ac.in}

\author{Manojkumar Patil}
\affiliation{%
  \institution{Indian Institute of Science}
  \city{Bengaluru}
  \country{India}}
\email{pmanojkumar@iisc.ac.in}

\author{Y Narahari}
\affiliation{%
  \institution{Indian Institute of Science}
  \city{Bengaluru}
  \country{India}}
\email{narahari@iisc.ac.in}

%%
%% By default, the full list of authors will be used in the page
%% headers. Often, this list is too long, and will overlap
%% other information printed in the page headers. This command allows
%% the author to define a more concise list
%% of authors' names for this purpose.
\renewcommand{\shortauthors}{Thandava, Geetha, Rohit, Priyanka, Manojkumar, and Narahari}

%%
%% The abstract is a short summary of the work to be presented in the
%% article.
\begin{abstract}

\end{abstract}

%%
%% The code below is generated by the tool at http://dl.acm.org/ccs.cfm.
%% Please copy and paste the code instead of the example below.
%%
\begin{CCSXML}
<ccs2012>
   <concept>
       <concept_id>10003752.10010070.10010099.10010102</concept_id>
       <concept_desc>Theory of computation~Solution concepts in game theory</concept_desc>
       <concept_significance>500</concept_significance>
       </concept>
   <concept>
       <concept_id>10002950.10003714.10003716</concept_id>
       <concept_desc>Mathematics of computing~Mathematical optimization</concept_desc>
       <concept_significance>500</concept_significance>
       </concept>
   <concept>
       <concept_id>10003752.10010070.10010099.10010100</concept_id>
       <concept_desc>Theory of computation~Algorithmic game theory</concept_desc>
       <concept_significance>500</concept_significance>
       </concept>
   <concept>
       <concept_id>10010405.10010476.10010480</concept_id>
       <concept_desc>Applied computing~Agriculture</concept_desc>
       <concept_significance>500</concept_significance>
       </concept>
 </ccs2012>
\end{CCSXML}
\ccsdesc[500]{Theory of computation~Algorithmic game theory}
\ccsdesc[500]{Theory of computation~Solution concepts in game theory}
\ccsdesc[500]{Mathematics of computing~Mathematical optimization}
\ccsdesc[500]{Applied computing~Agriculture}

%%
%% Keywords. The author(s) should pick words that accurately describe
%% the work being presented. Separate the keywords with commas.
\keywords{}
%% A "teaser" image appears between the author and affiliation
%% information and the body of the document, and typically spans the
%% page.
% \begin{teaserfigure}
%   \includegraphics[width=\textwidth]{sampleteaser}
%   \caption{Seattle Mariners at Spring Training, 2010.}
%   \Description{Enjoying the baseball game from the third-base
%   seats. Ichiro Suzuki preparing to bat.}
%   \label{fig:teaser}
% \end{teaserfigure}

\received{20 February 2007}
\received[revised]{12 March 2009}
\received[accepted]{5 June 2009}

%%
%% This command processes the author and affiliation and title
%% information and builds the first part of the formatted document.
\maketitle

\section{Introduction}
Agriculture forms the economic backbone of many developing countries, with the majority of farmers being small and marginal. In India, over 600 million people are dependent on farming for their livelihoods; the sector contributes approximately 17- 18\% of the country's gross domestic product and employs nearly half its workforce \cite{MOSPI2024}. At the heart of this vast agrarian landscape lies a defining structural characteristic: the overwhelming dominance of small and marginal farmers. According to the Agricultural Census, more than 89.4\% of farm holdings in India are smaller than two hectares, classifying their operators as small or marginal farmers \cite{MoAFW2023SmallFarmers}. These households cultivate fragmented plots, operate with limited capital, and remain acutely vulnerable to income shocks from erratic monsoons, volatile commodity prices, and rising input costs.

Against this backdrop, carbon farming has emerged as a promising practice. Carbon farming refers to a suite of land management activities, including conservation tillage, cover cropping, agro-forestry, improved irrigation, and biochar application, that enhance the capacity of agricultural soils to sequester atmospheric carbon dioxide. The scientific basis is well established: healthy, organically enriched soils act as a significant carbon sink, with global estimates suggesting that improved agricultural soil management could sequester several billion tonnes of \text{CO}$_2$ annually. Beyond their climate mitigation potential, these practices yield co-benefits of direct relevance to Indian farmers: improved soil organic matter, enhanced water retention, greater resilience to drought, and, in many cases, increased crop yields over the medium to long term. Carbon farming thus offers a rare alignment between environmental imperatives and agrarian welfare.

The translation of carbon sequestration into economic value occurs through carbon credit markets. Under voluntary carbon market frameworks, verified reductions or removals of greenhouse gas emissions can be certified and sold as carbon credits to emitters seeking to offset their footprint. For farmers, participation in such a programme would open an entirely new revenue stream: income derived not from what they grow, but from how they grow it. This additional income stream could benefit farmers in a meaningful way. Nevertheless, realising this potential is far from straightforward. Carbon credit programmes impose substantial fixed costs, including measurement, reporting, and verification (MRV) expenses, third-party auditing fees, and registry transaction charges, that are largely insensitive to farm size. For a smallholder cultivating one or two hectares, these costs can easily eclipse the value of credits generated, rendering solo participation economically non-viable.

This economic reality points directly to the importance of \textit{cooperative farming} as an enabling institution. By pooling their land and coordinating their carbon sequestration activities, groups of small farmers can collectively generate a credit volume that justifies shared certification costs, access markets that require minimum lot sizes, and present a credible counterparty to buyers and registries. Coalition formation among farmers is not a new idea in India; cooperative dairy federations, sugarcane cooperatives, and farmer producer organisations (FPOs) have demonstrated that collective action can substantially improve smallholder bargaining power and access to markets. 

Extending this cooperative logic to carbon farming, however, introduces a set of analytical challenges. Chief among these challenges is heterogeneity. Small and marginal farmers, even within the same village or watershed, differ substantially in farm size, soil composition, geographic conditions, and financial capacity. These differences imply that the carbon sequestration potential of any given practice, and the cost of adopting it, will vary across farmers. A practice that is highly profitable for one farmer may be infeasible or unprofitable for a neighbour. Consequently, farmers within a coalition may rationally choose different practice portfolios, introducing \textit{practice heterogeneity} that complicates both the measurement of coalition-level sequestration and the structure of shared certification costs. Furthermore, it is insufficient to ask merely whether a coalition is collectively beneficial; a viable and stable coalition must also distribute its monetary gains in a manner that every member prefers participation over independent action, which is a question of fair and individually rational payoff allocation.

In this work, we address these challenges through the lens of \textit{cooperative game theory}, a formal framework ideally suited to analysing the formation of coalitions and the distribution of collectively generated value among self-interested agents. Our specific contributions are as follows:

\begin{enumerate}
    \item \textbf{Optimal Carbon Practice Selection at the Individual Farmer Level.}
    We formulate the problem of selecting the most profitable combination of carbon sequestration practices for an individual farmer as a constrained optimisation problem. The objective accounts for carbon credit revenue, yield effects induced by each practice, and operational costs, subject to the farmer's budget constraint and compatibility restrictions between practices. Critically, the model incorporates pairwise synergistic interactions in both carbon sequestration and operational costs, capturing the reality that certain practice combinations amplify sequestration benefits or generate cost efficiencies beyond what each practice achieves independently. The solution yields each farmer's individually optimal practice portfolio, which serves as the foundation for coalition-level analysis.

    \item \textbf{Coalition Formation Efficiency and Monetary Distribution under Heterogeneous Practice Portfolios.}
    We study the cooperative game induced when small and marginal farmers form coalitions around carbon credit programmes, allowing each farmer to retain their individually optimal practice portfolio rather than imposing uniform practice adoption across the coalition. We define a characteristic function that captures coalition-level carbon credit revenue and shared certification costs, which scale with both total land area and the diversity of practices adopted, and analyse when coalition formation generates surplus value relative to independent participation. We then examine the distribution of this surplus, using cooperative game-theoretic solution concepts to characterise fair and stable payoff allocations that satisfy individual rationality for every coalition member. 
    % This formulation reveals a fundamental \textit{diversity--efficiency tradeoff}: heterogeneous practice portfolios allow farmers to optimise sequestration to their local soil and geographic conditions, but expand the set of distinct practices that must be verified, raising shared certification costs.
\end{enumerate}

The remainder of this paper is organised as follows. Section~2 reviews related work on carbon markets, cooperative approaches to smallholder agriculture, and the application of cooperative game theory in agricultural and environmental settings. Section~3 presents the formal methodology, beginning with the individual practice selection problem and proceeding to the coalition formation model. Section~4 describes the data. Section~5 presents experimental results, and Section~6 concludes with policy implications and directions for future research.

\section{Literature Review}
Existing research demonstrates the viability of farmer cooperatives for various agricultural applications. Bhardwaj et al. \cite{bhardwaj2023designing} leveraged farmer collectives to design fair and cost-optimal auctions for procuring agri-inputs. Sarkar et al. \cite{sarkar2023coalition} studied agricultural cooperatives using multi-agent coalitional models, showing positive revenue generation for smallholder farmers through similarity-based coalition formation. Zheng et al. \cite{zheng2022cooperative} developed cooperative game-based incentive mechanisms for biopesticide adoption through farmer-producer organizations (FPOs).

Recent work also explores cooperatives in carbon farming contexts. Mantey et al. \cite{mantey2025carbon} examined the role of dairy cooperatives in smallholder agricultural carbon projects, highlighting their importance in promoting sustainable farming. Wang et al. \cite{wang2024effects} investigated how farmer cooperatives reduce agricultural carbon emissions through reduced fertilizer inputs. Zhang et al. \cite{zhang2025can} systematically studied agricultural cooperatives’ impact on planting-industry carbon emissions using econometric models. Zheng et al. \cite{zheng2023economic} tested the role of marketing cooperatives in enhancing the economic performance of green agriculture.

Li et al. \cite{li2013carbon} evaluated farming practices based on carbon sequestration capacity and household income contribution in paddy fields. Piscitelli et al. \cite{piscitelli2023carbon} developed bottom-up approaches for identifying sustainable farming strategies to enhance soil carbon stocks. Mallappa et al. \cite{MALLAPPA2025100317} investigated farmer perceptions of carbon credit program actors and technology adoption motivations.

\section{Methodology}

Let $F = \{F_1, F_2, \ldots, F_n\}$ be the set of $n$ farmers willing to
participate in the carbon credit program. Each farmer $F_i$ is characterized
by their farm size $\text{FS}_i$ (acres). Let $P = \{P_1, P_2, \ldots, P_m\}$
denote the set of available carbon sequestration practices, where each practice
$P_j$ is associated with an operational cost $\text{OC}_j$ (rupees/acre/year)
and a carbon sequestration potential $\text{CSP}_j$ (tons~CO$_2$/acre/year).
We additionally define $\text{CSP}_{ij}$ as the realised carbon sequestration
potential of farmer $F_i$'s land under practice $P_j$, which may vary across
farmers due to differences in soil composition, geography, or local conditions.

\subsection{Optimal Practice Selection}

A farmer may simultaneously adopt a subset $\Pi \subseteq P$ of practices.
Since practices can interact among themselves, either amplifying sequestration
benefits or producing cost synergies and redundancies, we model these pairwise
interactions explicitly. Let $\alpha_{jk}$ denote the synergistic effect on
carbon sequestration when practices $P_j$ and $P_k$ are jointly adopted, and
let $\gamma_{jk}$ denote the corresponding pairwise cost interaction. The total
carbon sequestration potential and operational cost of a practice set $\Pi$ are
then defined as:

\begin{equation}
    \text{CSP}(\Pi) = \sum_{j:\, P_j \in \Pi} \text{CSP}_j
    + \sum_{\substack{P_j,\, P_k \in \Pi \\ j < k}} \alpha_{jk}
\end{equation}

\begin{equation}
    \text{OC}(\Pi) = \sum_{j:\, P_j \in \Pi} \text{OC}_j
    + \sum_{\substack{P_j,\, P_k \in \Pi \\ j < k}} \gamma_{jk}
\end{equation}

The farmer's objective is to select the practice set $\Pi^*$ that maximises
net profit comprising carbon credit revenue and yield-adjusted income, net of
operational costs, subject to a budget constraint and practice compatibility
constraints:

\begin{equation}
    \Pi^* = \arg\max_{\Pi \,\subseteq\, P}
    \Bigl[\,\text{CSP}(\Pi) \times \text{CCP} + Y(\Pi) - \text{OC}(\Pi)\,\Bigr]
\end{equation}
\begin{equation*}
    \text{s.t.} \quad \text{OC}(\Pi) \leq B
\end{equation*}

\noindent where the terms are defined as follows:
\begin{itemize}
    \item $\text{CSP}(\Pi)$: total carbon sequestration potential of practice set $\Pi$ (tons~CO$_2$/acre/year), as defined in Equation~(1).
    \item $\text{OC}(\Pi)$: total operational cost of implementing practice set $\Pi$ (rupees/acre/year).
    \item $\text{CCP}$: market price of one carbon credit (rupees/ton~CO$_2$).
    \item $Y(\Pi)$: incremental revenue attributable to yield changes induced by the adopted practices relative to baseline (rupees/acre/year).
    \item $B$: the farmer's maximum willingness-to-spend (rupees/acre/year).
\end{itemize}

The solution $\Pi^*$ identifies the most profitable combination of sequestration
practices feasible within the farmer's budget and compatible with the practice
constraints. We assume each farmer adopts their optimal portfolio $\Pi_i^*$
both when operating independently and when participating in a coalition; that
is, coalition membership does not alter individual practice selection.

\subsection{Coalition Formation with Heterogeneous Practice Portfolios}

Let each farmer $F_i \in S$ enter the coalition having independently selected
their optimal practice portfolio $\Pi_i^* \subseteq P$ as derived in
Section~3.1. Unlike a homogeneous setting, farmers within a coalition need not
adopt identical practices. Moreover, even when two farmers adopt the same
practice $P_j$, their realised sequestration $\text{CSP}_{ij}$ may differ due
to variation in soil composition, geography, or local conditions.

\subsubsection*{Farmer-Level Sequestration and Cost}

Given farmer $F_i$'s chosen portfolio $\Pi_i^*$, their individual realised
sequestration and operational cost are:

\begin{equation}
    \text{CSP}_i(\Pi_i^*) = \sum_{j:\, P_j \in \Pi_i^*} \text{CSP}_{ij}
    + \sum_{\substack{P_j,\, P_k \in \Pi_i^* \\ j < k}} \alpha_{jk}
\end{equation}

\begin{equation}
    \text{OC}_i(\Pi_i^*) = \sum_{j:\, P_j \in \Pi_i^*} \text{OC}_j
    + \sum_{\substack{P_j,\, P_k \in \Pi_i^* \\ j < k}} \gamma_{jk}
\end{equation}

Here $\text{CSP}_{ij}$ is farmer $F_i$'s realised sequestration under practice
$P_j$. Note the distinction from Equation~(1): whereas $\text{CSP}_j$ denotes
the practice-level baseline sequestration potential used in the individual
optimisation problem, $\text{CSP}_{ij}$ reflects farmer-specific soil and
geographic conditions. Consequently, even if $\Pi_i^* = \Pi_\ell^*$ for two
farmers $F_i$ and $F_\ell$, it generally holds that
$\text{CSP}_i(\Pi_i^*) \neq \text{CSP}_\ell(\Pi_\ell^*)$. The pairwise
interaction terms $\alpha_{jk}$ and $\gamma_{jk}$ remain practice-specific and
uniform across farmers.

\subsubsection*{Coalition Practice Profile and Aggregate Sequestration}

We define the \textbf{coalition practice profile} as the tuple
$\boldsymbol{\Pi}_S = \{\Pi_i^*\}_{i \in S}$. The coalition's total
sequestration aggregates farmer-specific contributions weighted by farm size:

\begin{equation}
    \text{CSP}(S, \boldsymbol{\Pi}_S) = \sum_{i \in S} \text{FS}_i \times \text{CSP}_i(\Pi_i^*)
\end{equation}

\subsubsection*{Certification Costs under Practice Heterogeneity}

When farmers adopt differing portfolios, the auditing burden scales not only
with total land area but also with the number of distinct practices that must
be verified across the coalition. Let $\mathcal{U}_S = \bigcup_{i \in S}
\Pi_i^*$ denotes the set of all distinct practices adopted within coalition $S$.
The three components of certification cost are:

\begin{equation}
    C_{\text{MR}}(S, \boldsymbol{\Pi}_S) = \text{Fixed}_{\text{MR}}
    + \text{Variable}_{\text{MR}} \times \left(\sum_{i \in S} \text{FS}_i\right)^{\!\alpha}
    + \phi_{\text{MR}} \times |\mathcal{U}_S|
\end{equation}

\begin{equation}
    C_{\text{V}}(S, \boldsymbol{\Pi}_S) = \text{Fixed}_{\text{V}}
    + \text{Variable}_{\text{V}} \times \left(\sum_{i \in S} \text{FS}_i\right)^{\!\beta}
    + \phi_{\text{V}} \times |\mathcal{U}_S|
\end{equation}

\begin{equation}
    C_{\text{T}}(S) = \text{Fixed}_{\text{T}} + \text{Variable}_{\text{T}} \times |S|
\end{equation}

where $\phi_{\text{MR}},\, \phi_{\text{V}} \geq 0$ are per-practice overhead
parameters capturing the additional methodological complexity introduced by
heterogeneous portfolios, and $\alpha$, $\beta$ are power-law scaling
parameters governing economies of scale in monitoring/reporting and
verification costs, respectively. This creates a natural
\textbf{diversity--efficiency tradeoff}: practice heterogeneity allows each
farmer to optimise sequestration to their land's specific conditions, but
raises shared certification costs as $|\mathcal{U}_S|$ grows.

\subsubsection*{Characteristic Function}

The characteristic function under heterogeneous practice portfolios is:

\begin{equation}
\begin{split}
v(S,\, \boldsymbol{\Pi}_S) = \;
    & \text{CSP}(S, \boldsymbol{\Pi}_S) \times \text{CCP} - [ C_{\text{MR}}(S, \boldsymbol{\Pi}_S) + C_{\text{V}}(S, \boldsymbol{\Pi}_S) + C_{\text{T}}(S)]
\end{split}
\end{equation}


\begin{equation}
\begin{split}
v(S,\, \boldsymbol{\Pi}_S) = \;
    & \left[\sum_{i \in S} \text{FS}_i \times \text{CSP}_i(\Pi_i^*)\right] \times \text{CCP} \\
    & - \left[\text{Fixed}_{\text{MR}} + \text{Variable}_{\text{MR}}
        \times \left(\sum_{i \in S} \text{FS}_i\right)^{\!\alpha}
        + \phi_{\text{MR}} \times |\mathcal{U}_S|\right] \\
    & - \left[\text{Fixed}_{\text{V}} + \text{Variable}_{\text{V}}
        \times \left(\sum_{i \in S} \text{FS}_i\right)^{\!\beta}
        + \phi_{\text{V}} \times |\mathcal{U}_S|\right] \\
    & - \left[\text{Fixed}_{\text{T}} + \text{Variable}_{\text{T}} \times |S|\right]
\end{split}
\end{equation}

% The first term aggregates each farmer's gross carbon credit revenue weighted by
% farm size and their portfolio-specific realised sequestration. The remaining
% terms reflect the shared certification costs borne collectively by the
% coalition. Operational costs $\text{OC}_i(\Pi_i^*)$ are farmer-specific and
% private, and are therefore handled through the individual rationality constraint
% rather than included in the coalition value.

\subsubsection*{Individual Rationality}

The singleton value for farmer $F_i$ under their individually optimal portfolio
$\Pi_i^*$ is:

\begin{equation}
\begin{split}
\tilde{v}\!\left(\{i\},\, \{\Pi_i^*\}\right) = \;
    & \text{FS}_i \times \text{CSP}_i(\Pi_i^*) \times \text{CCP}
      - \text{FS}_i \times \text{OC}_i(\Pi_i^*) \\
    & - \left[\text{Fixed}_{\text{MR}} + \text{Variable}_{\text{MR}} \times \text{FS}_i^{\,\alpha}
        + \phi_{\text{MR}} \times |\Pi_i^*|\right] \\
    & - \left[\text{Fixed}_{\text{V}} + \text{Variable}_{\text{V}} \times \text{FS}_i^{\,\beta}
        + \phi_{\text{V}} \times |\Pi_i^*|\right] \\
    & - \left[\text{Fixed}_{\text{T}} + \text{Variable}_{\text{T}}\right]
\end{split}
\end{equation}

\textit{Note that while $\text{OC}$ appears in $\tilde{v}(\{i\})$, it is
excluded from $v(S)$ because operational costs are private; the IR constraint
implicitly ensures each allocated payoff $x_i$ covers farmer $F_i$'s OC.}

Here $|\Pi_i^*|$ is the number of distinct practices adopted by farmer $F_i$
as a singleton, so $|\mathcal{U}_{\{i\}}| = |\Pi_i^*|$, and
$C_{\text{T}}(\{i\}) = \text{Fixed}_{\text{T}} + \text{Variable}_{\text{T}}$
since $|\{i\}| = 1$. A farmer who adopts a more diverse portfolio faces higher
solo certification costs, which raises their individual rationality threshold
and makes coalition participation harder to satisfy -- a subtle but
consequential modelling feature.

Since each farmer adopts the same optimal portfolio $\Pi_i^*$ regardless of
whether they operate independently or within a coalition, farmer $F_i$
participates in coalition $S$ only if their allocated payoff satisfies:

\begin{equation}
    x_i \geq \tilde{v}\!\left(\{i\},\, \{\Pi_i^*\}\right), \quad \forall\, i \in S
\end{equation}

The coalition is collectively rational if $\sum_{i \in S} x_i =
v(S, \boldsymbol{\Pi}_S)$ and individual rationality holds for all $i \in S$.

\subsubsection*{Special Case: Homogeneous Practice Adoption}

A natural benchmark arises when all farmers in the coalition adopt the same
practice portfolio, i.e.\ $\Pi_i^* = \Pi^*$ for all $i \in S$. In this
case the union $\mathcal{U}_S = \Pi^*$ is fixed regardless of coalition size,
so the per-practice overhead terms $\phi_{\text{MR}} \times |\mathcal{U}_S|$
and $\phi_{\text{V}} \times |\mathcal{U}_S|$ become constants independent of
$|S|$. The characteristic function reduces to:

\begin{equation}
\begin{split}
v(S,\, \Pi^*) = \;
    & \left[\sum_{i \in S} \text{FS}_i \times \text{CSP}_i(\Pi^*)\right] \times \text{CCP} \\
    & - \left[\text{Fixed}_{\text{MR}} + \text{Variable}_{\text{MR}}
        \times \left(\sum_{i \in S} \text{FS}_i\right)^{\!\alpha}
        + \phi_{\text{MR}} \times |\Pi^*|\right] \\
    & - \left[\text{Fixed}_{\text{V}} + \text{Variable}_{\text{V}}
        \times \left(\sum_{i \in S} \text{FS}_i\right)^{\!\beta}
        + \phi_{\text{V}} \times |\Pi^*|\right] \\
    & - \left[\text{Fixed}_{\text{T}} + \text{Variable}_{\text{T}} \times |S|\right]
\end{split}
\end{equation}

Note that farmer-specific sequestration terms $\text{CSP}_i(\Pi^*)$ still
permit heterogeneous contributions across farmers even under identical practice
adoption, since realised sequestration depends on individual soil and
geographic conditions. Comparing this homogeneous benchmark against the general
heterogeneous formulation isolates the net effect of practice diversity on
coalition value: the gain from allowing each farmer to optimise for their local
conditions on the revenue side, weighed against the increase in shared
certification costs from a larger $|\mathcal{U}_S|$.



\section{Datasets}
\section{Experimental Results}
\section{Conclusion and Future Work}


%% The acknowledgments section is defined using the "acks" environment
%% (and NOT an unnumbered section). This ensures the proper
%% identification of the section in the article metadata, and the
%% consistent spelling of the heading.
\begin{acks}

\end{acks}

%%
%% The next two lines define the bibliography style to be used, and
%% the bibliography file.
\bibliographystyle{ACM-Reference-Format}
\bibliography{sample-base}


%%
%% If your work has an appendix, this is the place to put it.
% \appendix

% \section{Research Methods}

% \section{Online Resources}

\end{document}
\endinput
%%
%% End of file `sample-sigconf.tex'.
